\chapter{Causal discovery as constrained scientific search}\label{ch:discovery}
\section{What a graph represents}
A bank may want to understand whether deposit-rate changes precede withdrawals,
whether a service complaint mediates attrition, or whether an economic shock
affects several products. A discovery procedure can help organize hypotheses,
but its output must state the variables, time scale, sampling mechanism,
latent-variable assumptions and edge semantics.

A DAG specifies acyclic causal relations among variables. A CPDAG represents
a Markov equivalence class under an appropriate causally sufficient model.
A PAG represents invariant ancestral information across a class that can
include hidden confounding and selection. Circle endpoints encode uncertainty;
they are not weak versions of directed arrows. Causal graphs also differ
from data-lineage graphs, whose edges represent transformations or provenance.

Conditional independences alone often leave multiple DAGs compatible with
the same observed law \citep{peters2017}. Temporal precedence rules out
some directions, but contemporaneous dependence, subsampling, delayed
recording and hidden common causes still matter. A graph of monthly
aggregates need not preserve daily intervention semantics.

\section{Comparison of method families}
\begin{longtable}{P{27mm}P{42mm}P{38mm}P{37mm}}
\toprule Method family & Identifying structure & Output/use & Failure to examine\\\midrule
\endhead
PC/FCI & Markov and faithfulness assumptions; suitable independence tests; different latent-variable assumptions & CPDAG or PAG; background constraints & Selection, measurement error, weak signals, finite-sample CI errors\\
LiNGAM / structural VAR & Linear equations and independent non-Gaussian innovations; acyclic contemporaneous structure & Oriented coefficients under the model & Gaussianity, hidden confounding, aggregation\\
NOTEARS / DAGMA & Smooth acyclicity constraint plus a chosen score and structural model & Numerical DAG candidate & Nonconvex solutions, sufficiency, thresholding, score equivalence\\
PCMCI & Stationary lagged system with suitable parent conditioning and tests & Lagged temporal relations under assumptions & Unmodeled contemporaneous dependence and omitted lags\\
PCMCI+ & Sufficiency, stationarity, adequate lag, Markov/faithfulness, valid tests & Lagged graph and contemporaneous equivalence information & Selection and measurement error excluded by the theorem\\
LPCMCI & Latent-confounded temporal ancestral-graph assumptions & Temporal PAG & Lag-window marginalization and unresolved endpoints\\
DYNOTEARS & Structural VAR with independent errors and functional/noise restrictions & Contemporaneous and lag coefficients & Fixed graph, stationarity/model fit, latent confounding\\
\bottomrule
\caption{Method assumptions determine what the output may assert. Runtime and finite-sample accuracy need separate evaluation.}
\label{tab:discovery}
\end{longtable}

The PC literature separates pointwise consistency from stronger
high-dimensional uniform consistency requirements \citep{kalisch2007}.
LiNGAM and its structural-VAR extension obtain additional directionality
through non-Gaussian independent noise \citep{shimizu2006,hyvarinen2010}.
Granger-style prediction instead asks whether past values improve
conditional prediction. It becomes interventionally causal only with
additional structural assumptions; a common macro driver can create
predictive lead-lag relations without a direct causal effect.

\section{Continuous optimization is a computational result}
NOTEARS replaces a combinatorial DAG constraint with
\begin{equation}
 h(W)=\operatorname{tr}\{\exp(W\circ W)\}-p=0.
 \label{eq:notears}
\end{equation}
For $S=W\circ W\geq0$, each $\operatorname{tr}(S^k)$ sums weights
of closed directed walks. Expanding the matrix exponential shows that
$h(W)=0$ precisely when no directed cycle exists
\citep{notears2018}. This elegant characterization makes continuous
optimization possible; it does not establish which direction is causal.
DAGMA gives another smooth characterization using a log determinant
\citep{dagma2022}. Both still need a suitable structural model, regularization
and numerical checks.

DYNOTEARS adds lagged coefficients to a structural VAR and constrains the
contemporaneous graph \citep{pamfil2020}. Time order already constrains
lagged directions, while identification of within-period effects depends
on functional and innovation assumptions. A group-wide monthly model that
mixes accounting identities, price-setting and endogenous demand is not
automatically within those assumptions.

\section{Time-series discovery and latent causes}
PCMCI reduces irrelevant conditioning by estimating useful parent sets
before momentary conditional-independence tests \citep{runge2019}.
PCMCI+ adds contemporaneous links. Its inspected assumption statement
requires causal sufficiency, stationarity, sufficient maximum lag,
appropriate faithfulness and consistent tests, and excludes selection and
measurement error \citep{runge2020}. Banking measurement defects therefore
need explicit treatment before applying its theorem.

LPCMCI permits hidden common causes through temporal ancestral-graph
representations \citep{gerhardus2020}. The result is a PAG with uncertainty,
not a fully observed latent-state reconstruction. Changing the lag window
can change the marginal graph being represented. Store the window,
conditioning tests, significance thresholds, multiplicity treatment,
background constraints and bootstrap stability with every result.

Missing data adds another graph. Conditioning on observed records can create
selection dependencies; naive imputation can create apparently strong edges
\citep{tu2019}. A null result can reflect weak measurements or low power,
not absence of an economic mechanism. The discovery workspace must keep
measurement and missingness assumptions visible.

\section{Continuous time, feedback and recent directions}
Dynamic SCMs model variables as paths and define interventions at that level
\citep{boeken2024}. Connections between dynamical systems, interventions and
equilibria require solvability and stability assumptions
\citep{bongers2022}. Recent causal SDE results extend Markov properties and
do-calculus under specified solvability conditions \citep{sde2026}.
These are relevant to irregular transactions and continuous-rate dynamics,
but do not yet provide an automatically qualified bank estimator.

The recent search found several useful research comparators.
ITSY addresses irregular series but explicitly assumes causal sufficiency,
stationarity and MCAR \citep{itsy2026}. LUCID routes corrections according
to pervasive or local confounding signatures \citep{lucid2026}.
The Causal-TS paper describes nonstationary temporal workflows and accelerated
independence tests \citep{causalts2026}. Their author-reported performance
has not been reproduced here; the most recent versions are preprints.
They motivate benchmarks rather than immediate default dependencies.

Group-level causal background knowledge can reduce a large search problem
\citep{vargas2026}. For example, an agreed temporal ordering of eligibility,
assignment, receipt and outcomes excludes implausible edges. An enterprise
ontology alone supplies names and relationships, not that causal ordering.
Domain experts must justify it.

\section{A discovery-to-study workflow}
Start with a narrowly defined question and a time-respecting dataset.
Record domain constraints, plausible omitted variables and observation
mechanisms. Run transparent baseline methods and selected challengers;
compare equivalence classes, graph stability and sensitivity to aggregation,
lag windows and missingness.

Next, choose a small set of economically meaningful candidate mechanisms.
For each one, state a proposed intervention and ask whether its effect is
identified in every graph still under consideration. If not, report a range
or identify the experiment or measurement that would distinguish models.
Only then create an effect-study specification with a documented graph
choice. The change from discovered hypothesis to accepted identification
assumption must be an explicit event in the evidence record.

This workflow can reduce wasted study design and suggest better measurement.
It should not promise to discover the bank's unique true causal graph from
all data, especially when policies adapt and hidden information drives actions.

\section{A financial graph needs a unit and an intervention}
Consider daily balances, the posted deposit rate, customer withdrawals,
funding transfers and a market rate. ``Rate causes balances'' is incomplete.
The rate may be an account-specific offer, a regional posted rate or a
group-average accounting yield. Balances may be end-of-day stocks, intraday
liquidity or monthly averages. A causal edge between two aggregate series
does not necessarily describe the effect of changing one customer's offer.

The graph specification should therefore name the unit, aggregation window,
decision actor and hypothetical intervention for every actionable node.
Some nodes are accounting identities. End balance equals beginning balance
plus net flows and valuation adjustments. Treating all terms as independent
stochastic variables can create deterministic conditional relations that
violate the generic faithfulness assumptions of a discovery algorithm.
The data layer should record the identity and choose a scientifically
appropriate variable representation.

The distinction is useful even when the final graph remains uncertain.
An economist may discover that all plausible graphs leave a rate effect
unidentified without a funding-cost instrument or an experiment. That is
an actionable output: obtain the missing variation or narrow the claim.
A unique-looking graph created by an arbitrary tie-breaker would be less
useful because it conceals the unresolved measurement decision.

\section{Why observational equivalence survives unlimited data}
For three variables, the DAGs $X\to M\to Y$ and
$X\leftarrow M\to Y$ have the same skeleton and no unshielded collider.
Both can imply $X\perp Y\mid M$ while allowing marginal dependence.
Under ordinary observational conditional-independence assumptions, the
data cannot distinguish their causal direction. A collider
$X\to M\leftarrow Y$ instead predicts marginal independence and
conditional dependence when its other assumptions hold. This is why
constraint-based discovery records equivalence information
\citep{peters2017,kalisch2007}.

The economic consequences differ. In the chain, intervening on $X$ can
change $Y$ through $M$. In the common-cause graph, an intervention on
$X$ may leave $Y$ unchanged. More observations can estimate the common
observed law more precisely without resolving that causal ambiguity.
Temporal order or an intervention can add information; a more confident
conditional-independence test alone cannot.

A similar ambiguity occurs in the executed two-variable VAR. Its lag
matrix and innovation covariance are
\[
 A=\begin{pmatrix}.7&0\\.15&.6\end{pmatrix},\qquad
 \Sigma_u=\begin{pmatrix}1&.8\\.8&1\end{pmatrix}.
\]
Ordering funding first gives impact column $(1,.8)'$ for a unit funding
shock. Ordering it second gives normalized column $(1,0)'$.
Both decompositions reproduce exactly the same Gaussian reduced-form law.
Their loan responses are the second components of $A^h(1,.8)'$ and
$A^h(1,0)'$, plotted in Figure~\ref{fig:dependence-irf}. The impact
response is $0.8$ in one model and zero in the other. No forecasting
comparison can choose between them under that common law.

Independent non-Gaussian structural innovations can break some linear
equivalences under the assumptions of LiNGAM and its time-series extension
\citep{shimizu2006,hyvarinen2010}. This is additional model information,
not a consequence of fitting a VAR. Hidden common causes, Gaussian
innovations, misspecified lag dynamics and aggregated signals can invalidate
the argument. The discovery interface should show which restriction
contributed each orientation.

\section{Conditional-independence testing is a statistical component}
For jointly Gaussian variables, zero partial correlation corresponds to
conditional independence. One can residualize $X$ and $Y$ on a
conditioning set $S$, calculate their correlation, and use an appropriate
sampling law. For nonlinear relationships, kernel, nearest-neighbor or
learned tests require different regularity and calibration. A strong
nonlinear dependence can coexist with zero ordinary correlation.

Temporal dependence complicates the sampling distribution. A test designed
for iid rows cannot be applied blindly to overlapping lag vectors. Parent
conditioning in PCMCI seeks to remove irrelevant persistence and reduce
the conditioning dimension before its final momentary tests
\citep{runge2019}. PCMCI+ incorporates contemporaneous structure under
its stated sufficiency and stationarity restrictions
\citep{runge2020}. These methods improve a specified search problem;
they do not exempt the observations from their own measurement model.

High-dimensional search also makes the testing family adaptive. With
$p$ variables and $L$ lags, even one test per candidate has order
$p^2L$ comparisons. Conditioning-set exploration can add many more.
An edgewise nominal level does not by itself imply graphwise error control.
Any false-discovery correction must match the tests and dependence it
actually receives. Bootstrap selection frequency is useful stability
information, but it is not automatically a posterior edge probability
or a calibrated confidence level.

The proposed benchmark therefore measures at least three objects.
First, test calibration on known independent or dependent processes under
the relevant serial dependence. Second, graph recovery under declared
causal models, reporting adjacency and orientation errors separately.
Third, the downstream effect or decision loss induced by using the graph.
A small structural Hamming distance can still conceal the one orientation
that reverses an important funding decision.

\section{Acyclic optimization and time-series structure}
In a linear contemporaneous model $X=XW+U$, a squared-error score plus
sparsity penalty and Equation~\eqref{eq:notears} defines an optimization
problem. The acyclicity condition is exact for the weighted adjacency
matrix, but identification also depends on what distribution of $U$ and
what score distinguish observationally equivalent parameterizations.
Thresholding small coefficients after optimization can alter the fitted
distribution and the uncertainty of the selected graph.

For lagged systems, write
\[
 X_t=B_0X_t+\sum_{\ell=1}^{L}B_\ell X_{t-\ell}+\varepsilon_t.
\]
Only $B_0$ needs a within-period acyclicity restriction in a finite
unrolled time graph. Solving gives reduced-form coefficients
$A_\ell=(I-B_0)^{-1}B_\ell$ and innovations
$u_t=(I-B_0)^{-1}\varepsilon_t$. Estimating $A_\ell$ alone
therefore does not identify the structural $B_\ell$. DYNOTEARS fits
contemporaneous and lagged structure together under its structural model
\citep{pamfil2020}; the product should retain both matrices rather than
present reduced-form predictive coefficients as direct mechanisms.

Feedback can be economically real within an observation window. Rates
affect flows and flows induce repricing during the same month. Forcing
an acyclic monthly contemporaneous graph may be a convenient approximation,
not a correct mechanism. Finer timing, a cyclic structural model with
well-defined solutions, or a path-level intervention can be more suitable.
Dynamic SCM and dynamical-systems work provides a language for these
questions \citep{boeken2024,bongers2022}, while practical estimation
still requires an observation scheme and identifying restrictions.

\section{Aggregation, missingness and changing regimes}
Sampling every second period of a VAR(1) produces
\[
 X_t=A^2X_{t-2}+Au_{t-1}+u_t.
\]
Even if $A$ is sparse, $A^2$ can contain a nonzero $(i,j)$ entry
because of an intermediate path $j\to k\to i$. The coarser lag graph
can therefore contain an edge that represents a multi-step mechanism.
Its innovation covariance also changes. Learning the monthly graph and
interpreting each edge as a daily direct intervention would be a category
error, not merely a loss of temporal precision.

Irregular observations add selection. An account can be measured more
often after an adverse event; staff may record a customer conversation
only when assistance is considered. Treating the resulting absence as
MCAR can create or remove conditional dependencies. Missingness-aware
discovery requires assumptions about the observation process
\citep{tu2019}. ITSY's explicit MCAR and stationarity restrictions
\citep{itsy2026} are useful precisely because they reveal when the method
does not fit a bank's informative observation schedule.

A pooled graph across low-rate, crisis and high-rate regimes can also mix
mechanisms. One approach fits separate graphs in sufficiently supported
regimes; another models changing parameters or environments under declared
invariance restrictions. Separate fits reduce sample size; a shared graph
can force false stability. The choice should be evaluated with temporal
holdouts and mechanism-specific diagnostics. A change-point detector
locates distributional change without establishing which causal mechanism
changed or whether the cause was policy.

Recent LUCID and Causal-TS proposals address important confounding or
nonstationarity challenges \citep{lucid2026,causalts2026}. They should be
included as research comparators on banking-specific synthetic mechanisms
and carefully bounded real studies. Their newness and author-reported
performance do not establish that they solve aggregation, hidden logging
or measurement selection simultaneously.

\section{From graph uncertainty to a useful answer}
Suppose a discovery method returns a set of compatible graphs
$\mathcal G$. For each $G$, let $\Theta_G$ be the compatible effect
set under its assumptions. The union $\bigcup_{G\in\mathcal G}\Theta_G$
describes what remains compatible when graph uncertainty is retained.
If a candidate policy is economically preferable across that set, the
decision can be robust even without unique orientation. If the set spans
both a gain and a loss, reporting one graph's point effect would hide a
material uncertainty.

The representation must distinguish an equivalence class justified by
theory from a convenient collection of bootstrap graphs. The latter is a
stability sample and need not contain every compatible graph. Likewise,
a PAG endpoint carries ancestral information under a latent-variable
model; replacing a circle by an arrow because its score is slightly larger
changes the identification claim.

A research workbench can then rank additional observations or interventions
by their potential to resolve the consequential ambiguity. For example,
randomized rate offers in a supported population may distinguish a
rate-response edge, while an independent liquidity measurement may help
separate demand from funding pressure. The first release can present these
as analyst-authored study proposals tied to graph alternatives. Automatic
optimal experiment selection would require a model of information value
and operating constraints and is a separate research feature.

\section{The discovery module's construction contract}
Inputs include variable definitions, population, observation clock,
aggregation, lag window, background constraints, hidden-variable and
selection assumptions, conditional-independence test and multiplicity
policy. Outputs include graph type, endpoint marks, edge evidence,
unresolved alternatives, stability information and sensitivity to timing
and missingness. Every display of a causal effect identifies the graph
and additional assumptions used to compute it.

The product should initially wrap established temporal and latent-aware
algorithms, retain their native diagnostics and compare simple structural
baselines. It should use domain groupings to reduce search only when those
groupings supply justified background restrictions
\citep{vargas2026}. Data-lineage links, semantic similarity and causal
arrows remain different edge types. Their integration in one interface
is useful only if their meanings remain explicit.
